\documentclass[a4paper, serif, 10pt]{moderncv}

%% moderncv
% color and style
\moderncvcolor{black}
\moderncvstyle{banking}

% renew moderncv command to adapt for CJK colon
\renewcommand*{\cvitem}[3][.25em]{%
  \ifstrempty{#2}{}{\hintstyle{#2}：}{#3}%
  \par\addvspace{#1}}

\renewcommand*{\cvitemwithcomment}[4][.25em]{%
  \savebox{\cvitemwithcommentmainbox}{\ifstrempty{#2}{}{\hintstyle{#2}：}#3}%
  \setlength{\cvitemwithcommentmainlength}{\widthof{\usebox{\cvitemwithcommentmainbox}}}%
  \setlength{\cvitemwithcommentcommentlength}{\maincolumnwidth-\separatorcolumnwidth-\cvitemwithcommentmainlength}%
  \begin{minipage}[t]{\cvitemwithcommentmainlength}\usebox{\cvitemwithcommentmainbox}\end{minipage}%
  \hfill% fill of \separatorcolumnwidth
  \begin{minipage}[t]{\cvitemwithcommentcommentlength}\raggedleft\small\itshape#4\end{minipage}%
  \par\addvspace{#1}}

%% page layout/margins
\usepackage[top=2.5cm, bottom=2.5cm, left=1.5cm, right=1.5cm]{geometry}

\nopagenumbers{}

%% Babel config for English language
\usepackage[english]{babel}

%% fontspec
\usepackage{fontspec}

\IfFontExistsTF{Linux Libertine}{
  \setmainfont[Ligatures={TeX, Common}, Numbers=Lining, ItalicFont=Linux Libertine]{Linux Libertine}
}{}
\IfFontExistsTF{Linux Libertine O}{
  \setmainfont[Ligatures={TeX, Common}, Numbers=Lining, ItalicFont=Linux Libertine O]{Linux Libertine O}
}{}

%% CTeX
% CJK support, used to show CJK characters in the resume
%
% - heading=false: disable ctex heading
% - punct=kaiming: use kaiming punctuations styles for CJK
% - scheme=plain: use plain scheme, do not override \normalsize font size
% - space=auto: space settings for CJK characters
%
% CJK fonts are resolved per locale language, e.g. Noto Serif/Sans CJK JP
% for Japanese resumes, with fallback to the next region if not installed
%
% ref:
% - http://ctan.mirrorcatalogs.com/language/chinese/ctex/ctex.pdf
\usepackage[UTF8, heading=false, punct=kaiming, scheme=plain, space=auto]{ctex}

\IfFontExistsTF{Noto Serif CJK SC}{
  \setCJKmainfont{Noto Serif CJK SC}
}{}
\IfFontExistsTF{Noto Sans CJK SC}{
  \setCJKsansfont{Noto Sans CJK SC}
}{}

%% line spacing
\usepackage{setspace}
\setstretch{1.125}

%% URL styling - use normal text instead of monospace
\urlstyle{same}

%% Patch \httplink and \httpslink to handle full URLs
% The original moderncv macros always prepend http:// or https:// to URLs.
% This causes issues when the URL already has a protocol (e.g., https://example.com)
% resulting in malformed URLs like https://https://example.com.
% This patch checks if the URL already contains "://" and uses it directly if so.
% Note: We use \str_set:Nx (x-expansion) to expand macros like \@homepage first.
\ExplSyntaxOn
\str_new:N \g_yamlresume_url_str

\RenewDocumentCommand{\httpslink}{O{}m}{
  \str_set:Nx \g_yamlresume_url_str {#2}
  \str_if_in:NnTF \g_yamlresume_url_str {://}
    { \url{#2} }
    { \url{https://#2} }
}

\RenewDocumentCommand{\httplink}{O{}m}{
  \str_set:Nx \g_yamlresume_url_str {#2}
  \str_if_in:NnTF \g_yamlresume_url_str {://}
    { \url{#2} }
    { \url{http://#2} }
}
\ExplSyntaxOff

\name{王雅婷}{}
\title{注册护士 | 心血管内科护理 | 8年临床经验}
\phone[mobile]{+86 138 1234 5678}
\email{wangyating.rn@outlook.com}
\homepage{https://www.linkedin.com/in/wangyating-rn}

\address{中国，北京市，北京，北京市朝阳区建国路88号院2号楼，100022}{}{}

\extrainfo{{\small \faLinkedin}\ \href{https://www.linkedin.com/in/wangyating-rn}{@wangyating-rn} {} {} {} • {} {} {} 
{\small \faWeixin}\ \href{https://weixin.qq.com/}{@WangYaTingRN}}

\begin{document}

\maketitle

\section{简介}

\cvline{}{\begin{itemize}
\item 注册护士，护理学学士，8年三级综合医院心血管内科及重症护理经验
\item 熟练掌握心电监护、静脉治疗、急救复苏、危重患者病情观察与护理文书书写
\item 具备护理组长经验，擅长排班协调、低年资护士带教及护理质量持续改进
\item 熟悉电子病历与护理信息系统，注重患者安全、健康宣教及多学科协作
\item 持有护士执业证书、BLS及ACLS证书，英语可阅读医学文献并进行日常沟通
\end{itemize}}
\section{教育背景}

\cventry{2012年9月–2016年6月}
        {学士，护理学，成绩：3.8/4.0}
        {北京协和医学院护理学院}
        {\url{https://www.pumc.edu.cn}}
        {}
        {\begin{itemize}
\item 系统学习临床护理、急危重症护理及护理管理核心课程
\item 在附属医院完成内科、外科、急诊及ICU轮转实习
\item 获校级护理技能竞赛二等奖，毕业成绩位列专业前15\%
\end{itemize}
\textbf{课程}：基础护理学、内科护理学、外科护理学、急危重症护理学、护理管理学、护理心理学、健康评估、护理伦理学}
\section{工作经历}

\cventry{2021年3月至今}
        {心血管内科护理组长}
        {北京安和医院}
        {\url{https://www.anhehospital.cn}}
        {}
        {\begin{itemize}
\item 负责心血管内科病区护理小组日常管理，协调床位、排班及危重患者护理资源
\item 参与制定并落实跌倒、压疮、静脉血栓等风险评估与闭环管理流程
\item 带教新入职护士及实习护生，完成护理操作示范与考核
\item 组织科室业务学习和应急演练，提升团队急救配合与突发事件处置能力
\item 参与护理质量指标监测与PDCA持续改进，推动护理文书规范化
\end{itemize}
\textbf{关键字}：护理管理、心血管护理、患者安全、临床带教、质量改进}

\cventry{2018年7月–2021年2月}
        {注册护士}
        {北京仁济心血管病医院}
        {\url{https://www.renjiheart.cn}}
        {}
        {\begin{itemize}
\item 负责心血管内科患者入院评估、生命体征监测、用药观察及健康宣教
\item 熟练执行心电监护、静脉穿刺、输液泵和微量泵使用及危重患者护理
\item 配合医生完成介入术后护理、胸痛中心流程及抢救工作
\item 规范书写护理记录，参与交接班和不良事件上报
\end{itemize}
\textbf{关键字}：心电监护、静脉治疗、介入护理、健康宣教、急救配合}

\cventry{2016年7月–2018年6月}
        {注册护士}
        {北京康泰综合医院}
        {\url{https://www.kangtai-hospital.cn}}
        {}
        {\begin{itemize}
\item 在内科病区完成基础护理、生命体征监测、给药及护理记录书写
\item 参与夜班轮值与患者安全巡查，熟悉住院患者全流程护理
\item 接受急救技能培训并参与科室应急演练
\end{itemize}
\textbf{关键字}：基础护理、内科护理、患者安全、护理记录}
\section{语言}

\cvline{中文}{母语或双语水平 \hfill \textbf{关键字}：普通话二级甲等}
\cvline{英语}{完全专业水平 \hfill \textbf{关键字}：CET-6、医学英语阅读}
\section{专业技能}

\cvline{临床护理}{专家 \hfill \textbf{关键字}：护理评估、静脉穿刺、心电监护、给药管理、伤口护理}
\cvline{急危重症护理}{高级 \hfill \textbf{关键字}：BLS、ACLS、除颤仪、抢救配合、病情观察}
\cvline{护理管理}{高级 \hfill \textbf{关键字}：排班协调、临床带教、质量改进、PDCA、不良事件分析}
\cvline{护理信息系统}{中级 \hfill \textbf{关键字}：电子病历、护理文书、数据统计、办公软件}
\section{荣誉}

\cventry{2023年5月}
        {北京安和医院}
        {年度优秀护士}
        {}
        {}
        {\begin{itemize}
\item 因在心血管内科临床护理、患者安全及团队协作方面表现突出获奖
\item 年度护理质量考核优秀，患者满意度位于科室前列
\end{itemize}}

\cventry{2015年11月}
        {北京协和医学院护理学院}
        {护理技能竞赛二等奖}
        {}
        {}
        {在无菌操作、静脉输液及心肺复苏项目中表现优异。}
\section{证书}

\cventry{2016年8月}
        {北京市卫生健康委员会}
        {中华人民共和国护士执业证书}
        {}
        {}
        {}

\cventry{2024年3月}
        {美国心脏协会}
        {基础生命支持（BLS）证书}
        {\url{https://www.heart.org/en/cpr}}
        {}
        {}

\cventry{2024年4月}
        {美国心脏协会}
        {高级心血管生命支持（ACLS）证书}
        {\url{https://www.heart.org/en/cpr}}
        {}
        {}

\cventry{2022年9月}
        {北京护理学会}
        {心血管专科护士培训证书}
        {}
        {}
        {}
\section{出版刊物}

\cventry{2021年6月}
        {心血管内科老年患者跌倒风险分级护理实践}
        {中华护理杂志}
        {\url{https://www.zhhlzzs.com}}
        {}
        {\begin{itemize}
\item 总结老年心血管疾病患者跌倒风险评估与分级护理经验
\item 通过多学科协作和健康教育降低住院期间跌倒发生率
\item 为病区护理质量改进提供实践参考
\end{itemize}}
\section{项目}

\cventry{2022年1月–2022年12月}
        {针对心血管内科住院患者跌倒风险开展评估、干预、追踪和复盘的护理质量改进项目。}
        {心血管内科跌倒风险闭环管理项目}
        {\url{https://www.anhehospital.cn/nursing-quality}}
        {}
        {\begin{itemize}
\item 建立入院、转科、出院前动态跌倒风险评估流程
\item 设计高风险患者标识、床旁宣教和环境改造清单
\item 通过月度数据分析和PDCA循环持续优化护理措施
\item 项目实施后病区跌倒发生率较上一年度下降
\end{itemize}
\textbf{关键字}：患者安全、跌倒预防、PDCA、护理质量}
\section{兴趣爱好}

\cvline{专业学习}{护理文献阅读、急救培训、循证护理}
\cvline{生活爱好}{游泳、瑜伽、旅行}
\section{志愿服务}

\cventry{2019年3月–2023年12月}
        {健康宣教志愿者}
        {朝阳区社区健康服务中心}
        {\url{https://www.cyqhealth.org}}
        {}
        {\begin{itemize}
\item 参与社区高血压、冠心病患者健康宣教和用药指导
\item 协助开展心肺复苏公众培训与义诊活动
\item 为老年居民提供血压监测、饮食及运动建议
\end{itemize}}
    
\end{document}